% !TEX root =../main.tex

\section{Conclusiones}
A partir de el desarrollo de la pactica:
Un Hub es un dispositivo de red que funciona
en la capa 1 de el modelo OSI que permite
a otros dispositivos de red conectarse entre 
sí al retransmitir los datos que llegan por 
un puerto Ethernet a todos los demás puertos, 
solo permite transmisión half-duplex, lo que significa
que reciben información o envían, pero no en simultáneo.
Pueden ser utilizados para conectar varios dispositivos 
de red cuando el ancho de banda no es exigente.\cite{hub_coursera}
Al ser un dispositivo de capa 1 no tiene ninguna
capacidad de configurarse por lo que no se le puede brindar seguridad.
\\
El Protocolo de control de mensajes de Internet (ICMP)
es un protocolo de capa 3 de el modelo OSI que se utiliza para 
diagnosticar problemas en la comunicación de la red.
Su uso principal permite determinar si los datos llegan o no 
en un tiempo dado. Una prueba ICMP puede 
realizarse con herramientas como traceroute y ping, de las cuales ping
prueba la velocidad de la conexión entre dos dispositivos y informa el 
tiempo que se tarda un paquete de datos en llegar al destino y volver a remitente.
El protocolo ICMP tiene su propio paquete que contiene un encabezado ICMP y un
encabezado de direccionamiento IP.\cite{icmp_cloudflare}
\\
Una conexión Telnet es una transmisión sobre TCP usada para transmitir datos
intercalados con información de control Telnet.\cite{telnet_RFC} El protocolo Telnet es utilizado
para el acceso remoto por red a computadores (o dispositivos de red), es un protocolo
de cliente-servidor que provee al usuario una sesión de terminal  al dispositivo 
remoto desde la aplicación Telnet del cliente. Este protocolo no ofrece ningún tipo
de seguridad por lo que es vulnerable a ataques en los que un agente pueda escuchar 
los datos transmitidos. Hoy en día se utiliza SSH(Secure Shell) como alternativa
segura para la conexión remota.\cite{telnet_shh_academy}
\\